\chapter{Barium Enema}

\section*{Overview}
A barium enema, also called a lower gastrointestinal (GI) series, is a fluoroscopic X-ray examination of the colon and rectum. Barium contrast outlines the inside of the large intestine so that its shape and lining can be evaluated.

\section*{Alternative Names}
Lower GI series, lower GI tract radiography, double-contrast barium enema

\section*{Principle}
Barium sulfate is introduced through the rectum and coats the colonic lining. X-ray images are obtained as the patient changes position; in a double-contrast study, air or carbon dioxide is also introduced to expand the colon and improve visualization of the mucosa.

\section*{History}
Barium examinations of the gastrointestinal tract became established after the introduction of radiopaque contrast in the early twentieth century. Modern fluoroscopy and double-contrast techniques provide more detailed images, although colonoscopy and CT-based studies are preferred for many indications.

\section*{Why It Is Done}
It may be ordered to evaluate changes in bowel habits, abdominal pain, rectal bleeding, suspected narrowing or obstruction, diverticular disease, inflammatory changes, or an abnormality that cannot be fully assessed by another study.

\section*{Is This a Primary or Secondary Test or Procedure}
It is a diagnostic imaging test and is generally a secondary or alternative examination when colonoscopy is incomplete, unsuitable, unavailable, or not required for the clinical question. It does not allow tissue biopsy or polyp removal.

\section*{How It Is Performed}
You lie on an X-ray table, usually on your left side. A lubricated enema tube is gently inserted into the rectum, and barium is allowed to flow into the colon. The radiologist or technologist takes fluoroscopic images while you change position. The tube is then removed and you may be asked to expel the contrast in a bathroom before additional images are obtained.

\section*{Preparation}
The colon must be empty. Follow the facility's instructions for a low-residue or clear-liquid diet, laxatives, and possibly cleansing enemas. Tell the radiology team if you may be pregnant, have severe abdominal pain, suspected bowel obstruction or perforation, recent bowel surgery, or difficulty controlling your bowel movements. Do not change prescribed medicines unless instructed.

\section*{Contraindications and Precautions}
Avoid a barium enema when bowel perforation is suspected or when there is severe acute colitis, toxic megacolon, or a high risk of perforation. A water-soluble iodinated contrast agent may be chosen instead when perforation is a concern. Pregnancy requires individualized assessment because the examination uses ionizing radiation.

\section*{Risks}
Temporary cramping, bloating, and an urgent need to evacuate are common. Rare complications include contrast impaction, bowel perforation, infection, or inflammation if barium escapes through an unrecognized hole. The examination also involves a small amount of ionizing radiation. Allergic reactions to barium are extremely uncommon.

\section*{Aftercare and Recovery}
You can usually return home soon after the examination. Drink fluids as advised and follow instructions about resuming food and medicines. Stools may be pale or white for a short time. Contact a clinician promptly for severe or persistent abdominal pain, marked swelling, fever, vomiting, dizziness, or inability to pass stool or gas.

\section*{Outcome and Follow-up}
The images are reviewed for the contour and caliber of the colon, the appearance of the mucosal lining, diverticula, strictures, obstruction, masses, or other abnormalities. An abnormal result may require colonoscopy, CT, biopsy, or another test for confirmation.

\section*{Expected Results}
The colon and rectum have normal caliber and position, with an intact, smooth mucosal outline and no evidence of mass, obstruction, significant inflammation, or diverticular abnormality.

\section*{Follow-up}
Follow-up depends on the findings and the reason for testing. Colonoscopy is commonly used when direct inspection or biopsy is needed; CT or other imaging may be recommended for suspected complications or disease outside the bowel wall.

\section*{When to Seek Medical Care}
Follow the procedure team's specific instructions. Seek urgent medical assessment for severe or worsening pain, uncontrolled bleeding, fever or signs of infection, trouble breathing, chest pain, fainting, new weakness or confusion, inability to urinate, or any rapidly worsening symptom. The appropriate warning signs depend on the procedure and the patient's medical condition.

\section*{References}
\begin{enumerate}
  \item MedlinePlus. Barium enema. U.S. National Library of Medicine; reviewed January 1, 2025.
  \item American College of Radiology and Radiological Society of North America. Lower GI X-ray (Barium Enema). RadiologyInfo.org; reviewed April 22, 2025.
\end{enumerate}

\clearpage
